\begin{titlepage}
\color{navy}\sffamily\raggedright
\vspace*{30pt}
{\large THE ENERGY BALANCE PROJECT\par}
\vspace{30pt}
{\fontsize{32}{38}\selectfont\bfseries Mathematical Structure and Proof\par}
\vspace{20pt}
{\Large Part II\par}
\vspace{12pt}
{\large Revised Gaussian Explanatory Edition\par}
\vfill
{\large Konrad Grzyb\par}
\vspace{12pt}
{\normalsize Author-review manuscript\par September 2026\par}
\end{titlepage}
\chapter*{About this revised edition}
\addcontentsline{toc}{chapter}{About this revised edition}
The first volume asks what a physical account should describe. This volume
derives the mathematical statements that such an account can honestly make.
Its active worked model is the Local Gaussian specialization: a fixed
quadratic deviation potential on a conserved, nonnegative scalar resource.
The potential supplies a valuation, not an automatic action selector.

This edition revises the original unified explanatory Part II. It preserves
its patient approach to notation, derivations, assumptions and worked
calculations. The original editable manuscript was not recovered; the active
chapters here are reconstructed editable source, not a facsimile. The supplied
original PDF remains unchanged. Selected original proofs appear at the end
as a clearly labelled historical appendix.

Chapter numbering follows the revised Book I: this volume begins at 33 and
ends at 46. Original chapter and equation numbers survive only inside the
historical appendix, where they identify the old model. A repository
disposition ledger maps every original outline entry to its revised subject
or preserved historical home. Do not confuse an old equation number with a
new one simply because both books use the word EBU.

\begin{cautionbox}{Four different kinds of statement}
A definition declares a model. A proof establishes a consequence of explicit
assumptions. An implementation carries out a specified operation. An
experiment supplies evidence about a specified system. None becomes another
merely by being written in mathematical notation or committed to a repository.
\end{cautionbox}

The Gaussian programme remains prospective. Its capacity rule, ownership map,
event profile and random selection semantics require their own acceptance.
The conditional derivations below do not adopt them by implication. This
book introduces no scientific run, no parameter calibration and no claim of
economic or ecological success. It is not independently peer reviewed.

\chapter*{How to read the proofs}
\addcontentsline{toc}{chapter}{How to read the proofs}
Read Chapters 33--35 if vectors, units or action graphs are unfamiliar. They
turn a physical story into declared objects without importing a controller.
Chapters 36--40 derive the Gaussian field, exact finite accounting, locality,
path integration and telescoping. Chapters 41--43 address groups, attribution
and the proposed capacity account. Chapters 44--45 explain why a correct
account is not yet a theorem about dynamics. Chapter 46 combines the ideas
in a worked laboratory, with deliberately adverse as well as favorable cases.

A displayed equality is worth reading in two passes. First identify what
every term denotes and which state it uses. Then check the algebra. Most
serious accounting mistakes in this subject are not failures to multiply:
they compare different baselines, mix rates and quantities, omit an affected
factor, or add an institutional meaning that the equality did not prove.

The examples use small, exact numbers so the reader can reconstruct them.
They are illustrative arithmetic, not sampled trajectories. When a later
chapter writes a candidate transition rule, that is a specification on paper;
it is not a report that the research engine was executed to produce it.

The historical appendix preserves the original continuous-time, timestep,
network-certificate and P1C derivations. Read its chapter numbers as archival
identifiers. Its flow law is not the random-affordable actor process described
in the active programme. The new Chapter 44 explains the bridge and its limits.
